% Iteration 5 change manifest as two structured cards: chg-7 introduces the
% prompt + tool-descriptor publish-state rules, chg-8 installs the stateful
% shell guard. Both ship together at the iteration-4 boundary as the
% iteration-5 harness. Palette and layout follow case-trajectory.tex.
\begin{figure}[t]
    \centering
    \scriptsize
    \begin{tcbraster}[raster columns=2,
                      raster equal height,
                      raster column skip=4pt,
                      raster row skip=0pt,
                      raster left skip=0pt,
                      raster right skip=0pt,
                      boxrule=0.6pt,
                      left=4pt,right=4pt,top=3pt,bottom=3pt,
                      coltitle=white]
        \begin{tcolorbox}[colback=aheSky!18,colframe=aheNavy,colbacktitle=aheNavy,
                          title={\scriptsize\textbf{chg-7，迭代 5，commit \texttt{3ba3a90}，层级：提示词 + 工具描述符}}]
            \textbf{\textcolor{aheNavy}{涉及文件}} \\
            \rule{\linewidth}{0.3pt}\\[1pt]
            $\bullet$ \texttt{workspace/systemprompt.md} \\
            $\bullet$ \texttt{tool\_descriptions/run\_shell\_command.tool.yaml} \\[3pt]
            \rule{\linewidth}{0.3pt}\\
            \textbf{\textcolor{aheNavy}{改动内容}} \\
            \rule{\linewidth}{0.3pt}\\[1pt]
            向 harness 的工作记忆追加了三条规则。发布态规则：一旦评测器式的最终检查通过，由此得到的文件系统与服务状态就是交付面，不得为了``看起来干净''而重置。临时目录规则：把探索性产物放在 \texttt{/tmp} 或验证器会忽略的临时路径之下。字面输出规则：对于带有字节级契约的 DSL、配置或脚本输出，要在字节层面校验相等性。 \\[3pt]
            \rule{\linewidth}{0.3pt}\\
            \textbf{\textcolor{aheNavy}{修复的失败模式}} \\
            \rule{\linewidth}{0.3pt}\\[1pt]
            智能体已经到达评测器可通过的状态，随后为了``收拾整洁''而发出大范围清理命令或重写输出，使验证器拿不到任何交付物。 \\[3pt]
            \rule{\linewidth}{0.3pt}\\
            \textbf{\textcolor{aheNavy}{预测修复的任务}} \\
            \rule{\linewidth}{0.3pt}\\[1pt]
            4 个任务。例如：\texttt{path-tracing}、\texttt{configure-git-webserver}、\texttt{polyglot-rust-c}、\texttt{large-scale-text-editing}。
        \end{tcolorbox}
        \begin{tcolorbox}[colback=aheTeal!12,colframe=aheBlue,colbacktitle=aheBlue,
                          title={\scriptsize\textbf{chg-8，迭代 5，commit \texttt{4e0aab9}，层级：工具实现}}]
            \textbf{\textcolor{aheBlue}{涉及文件}} \\
            \rule{\linewidth}{0.3pt}\\[1pt]
            $\bullet$ \texttt{tools/shell\_tools/run\_shell\_command.py} \\[3pt]
            \rule{\linewidth}{0.3pt}\\
            \textbf{\textcolor{aheBlue}{改动内容}} \\
            \rule{\linewidth}{0.3pt}\\[1pt]
            在 shell 工具内部安装了一个有状态的发布态守卫。在评测器式的最终检查成功之后，守卫会解析验收命令，从中提取显式的文件路径与根目录，并记录为受保护对象。此后若出现会删除受保护输出或重置受保护根目录的破坏性命令，将在执行前被拦截，并以针对性的报错返回。显式的 \texttt{ALLOW\_POST\_SUCCESS\_RESET} 令牌可以把该阻断降级为警告，此后智能体必须在提交前重新验证。 \\[3pt]
            \rule{\linewidth}{0.3pt}\\
            \textbf{\textcolor{aheBlue}{修复的失败模式}} \\
            \rule{\linewidth}{0.3pt}\\[1pt]
            即使已经有了提示词层面的规则，智能体在进入发布态之后仍会发出破坏性的清理命令。在 shell 工具处施加执行期强制是最直接的互锁手段。 \\[3pt]
            \rule{\linewidth}{0.3pt}\\
            \textbf{\textcolor{aheBlue}{预测修复的任务}} \\
            \rule{\linewidth}{0.3pt}\\[1pt]
            同样的 4 个任务；在 \texttt{path-tracing} 上起关键作用，该任务的 F4 步骤正是对 \texttt{/app/reconstructed.ppm} 执行 \texttt{rm -rf}。
        \end{tcolorbox}
    \end{tcbraster}
    \caption{在迭代 4 边界处一并写下、并作为迭代 5 的 harness 上线的两条变更清单条目。\texttt{chg-7} 在系统提示词与工具描述符中明确了发布态规则；\texttt{chg-8} 则在 shell 工具内部安装了执行期互锁。这一对变更在下一轮把 \texttt{path-tracing} 翻转为通过。}
    \label{fig:manifest-publish-state}
\end{figure}
